\documentclass{article}
\usepackage{amsthm, amssymb, latexsym, amsmath, amsbsy, graphicx, epsfig}
\usepackage{ mathrsfs }
\usepackage{ textcomp }

\title{Writing Assignment \#8: Complete Induction and Relations}
\author{Robert T Kohn}
\date{March 31st, 2025}

\begin{document}

\maketitle

{\huge{Section 2.5}}
\section*{Question 1}
\subsection*{Part B}

Claim: $(\forall \ n>22 \in N) (\exists s, t \in Z) \ \epsilon \ n=3s+4t \wedge s \geq 3 \wedge t \geq 2$

Case n=23: s=5, t=2

$3(5)+4(2)=15+8=23$

Case n=24: s=4, t=3

$3(4)+4(3)=12+12=24$

Case n=25: s=3, t=4

$3(3)+4(4)=9+16=25$

Then, $\exists k \in Z \epsilon 23 \leq k \leq n$ 

Thus,  $n \geq 25$

So, $n+1=(n-2)+3$

Hence, $n-2\geq  23 \Rightarrow n+1=3s+4t+3$

Then, $n+1=3(s+1)+4t$

$\therefore (\forall \ n>22 \in N) (\exists s, t \in Z) \ \epsilon \ n=3s+4t \wedge s \geq 3 \wedge t \geq 2$

\subsection*{Part C}

Claim: $(\forall \ n>33 \in N) (\exists s, t \in Z) \ \epsilon \ n=4s+5t \wedge s \geq 3 \wedge t \geq 2$

Case n=34: s=6, t=2

$4(6)+5(2)=24+10=34$

Case n=35: s=5, t=3

$4(5)+5(3)=20+15=35$

Case n=36: s=4, t=4

$4(4)+5(4)=16+20=36$

Case n=37: s=3, t=5

$4(3)+5(5)=12+25=37$

Then, $n+1=(n-4)+5$

Thus, $n-4 \geq 34 \Rightarrow n+1=4s+5t+5$

So, $n+1=4s+5t+5 \Rightarrow n+1=4s+5(t+1)$

$\therefore (\forall \ n>33 \in N) (\exists s, t \in Z) \ \epsilon \ n=4s+5t \wedge s \geq 3 \wedge t \geq 2$

\section*{Question 3}

Claim: $a_1=2 \wedge a_2=4 \wedge a_{n+2}=5a_{n+1}-6a_n \wedge n \geq 1 \Rightarrow a_n=2^n$

Case n=1:

$a_1=2^1=2$

Case n=2: 

$a_2=2^2=4$

Then, for some arbitrary k, $a_k=2^k$

Thus, $a_{n+1}=5*a_n-6a_{n-1}$

So, $a_{n+1}=5*2^n-6*2^{n-1}$

Hence, $a_{n+1}=2^{n-1}(5*2-6*1)=2^{n-1}*4=2^{n-1}*2^2$

Then, $a_{n+1}=2^{n+1}$

$\therefore a_1=2 \wedge a_2=4 \wedge a_{n+2}=5a_{n+1}-6a_n \wedge n \geq 1 \Rightarrow a_n=2^n$



\section*{Question 12}
\subsection*{Part A}

\

$g1=2; \ $

$g2=2; \ $

$g3=g2*g1=2*2=4 ; \ $

$g4=g3*g2=4*2=8; \ $

$g5=g4*g3=8*4=32;$

\subsection*{Part B}

Claim: $g_n=2^{f_n}$

Case n=1: 

$g_1=2^{f_1}=2$

Case n=2:

$g_2=2^{f_2}=2$

Then, for some arbitrary k, $g_i=2^{f_i}$ for all $i \in \{1,2...k\}$

Thus, $g_{k+1}=g_k*g_{k-1}$

So, $g_{k+1}=2^{f_k}*2^{f_{k-1}}=2^{f_k+f_{k-1}}$

Hence, $f_{k+1}=f_k+f_{k-1}$

$\therefore g_n=2^{f_n}$

\

\

{\huge{Section 3.1}}
\section*{Question 6}
\subsection*{Part K}

$R_{11}^{-1} = \{(x, y) \in P \times P \mid x \text{ is a sibling of } y\}$

\subsection*{Part L}

$R_{12}^{-1} = \{(x, y) \in P \times P \mid x \text{ loves } y\}$

\section*{Question 7}
\subsection*{Part C}

\

$(2,4) \in S \rightarrow (4,z) \in T \rightarrow (2,1)$

$(3,4) \in S \rightarrow (4,z) \in T \rightarrow (3,1) $

$(3,1) \in S \rightarrow (3,z) \in T \rightarrow (3,4) $

$(5,5) \in S \rightarrow (5,z) \in T \rightarrow (5,z) \notin T$

T \textopenbullet \  S = $\{(2,1),(3,1),(3,4)\}$


\# I don't know how much work you wanted to see here so I hope this is what you're looking for and let me know if this is too much as well

\subsection*{Part G}

\

$(1,4) \in T \rightarrow (4,z) \notin S$

$(3,5) \in T \rightarrow (5,z) \in S \rightarrow (3,5) $

$(4,1) \in T \rightarrow (1,z) \notin S $

S \textopenbullet \  T = $\{(3,5)\}$

$(3,5) \in $S \textopenbullet \  T $ \rightarrow (5,z) \in R \rightarrow (3,2) $

R \textopenbullet \  (S \textopenbullet \  T)= $\{(3,2)\}$

\


\section*{Question 11}
\subsection*{Part A}

Claim: Rng($R^{-1}$)=Dom(R)

Case 1: Rng($R^{-1}$) $\subseteq$ Dom(R)

Suppose $x \in Rng(R^{-1}$),

Then, $\exists y \epsilon (y,x) \in R^{-1}$

Thus, $(x,y) \in R$

So, $x \in Dom(R)$

Case 2: Dom(R) $\subseteq$ Rng($R^{-1}$)

Suppose $a \in Dom(R)$,

Then, $\exists b \epsilon (a,b) \in R$

Thus, $(b,a) \in R^{-1}$

So, $a \in Rng(R^{-1})$

$\therefore Rng(R^{-1})=Dom(R)$

\subsection*{Part B}

Claim: Dom(S \textopenbullet R) $\subseteq$ Dom(R)

Suppose x $\notin$ Dom(R)

Then, $\nexists y \epsilon (x,y) \in R$

Thus, $\nexists z (x,z) \in $ S \textopenbullet R)

Hence, $ x \notin Dom($S \textopenbullet R)

$\therefore$ Dom(S \textopenbullet R) $\subseteq$ Dom(R)

\subsection*{Part C}

Claim: Dom(S \textopenbullet R) = Dom(R)

Suppose $A=\{1\}, B=\{2,3\}, C=\{4\}, S=\{1,2\}, R=\{3,4\} $

S \textopenbullet R=$\emptyset \Rightarrow$ Dom(S \textopenbullet R)=\{$\emptyset$\}

Dom(R)=\{1\}

Dom(S \textopenbullet R) $\neq$ Dom(R)

$\therefore$ The inequality fails.

\section*{Question 14}

Claim: $A=m \wedge B=n \Rightarrow 2^{mn}$ different relations from A to B

Case m=1:

Then, $A=\{1\} \wedge B=n$

Thus, there are $2^n$ possible subsets.

Suppose $A=\{k\} \Rightarrow 2^{kn} \ relations$

Then, $A=\{k+1\} \Rightarrow 2^{(k+1)*n}$

$\therefore A=m \wedge B=n \Rightarrow 2^{mn} $ different relations from A to B

\section*{Question 17}
\subsection*{Part A}

This proof gets a grade of F. The claim is false and it assumes that it's composition and inversion commute. 

\subsection*{Part B}

This proof gets an A.

\subsection*{Part F}
This proof gets a grade of A.

\end{document}
